\documentclass[10pt,a4paper]{amsart}
\usepackage[margin=19mm]{geometry}
\usepackage{amsmath,amssymb,mathtools}
\usepackage[scr=rsfs,cal=euler]{mathalfa}
\usepackage{xfrac}
\usepackage[unicode,hidelinks,bookmarks=false]{hyperref}
\newcommand{\ie}{i.e.\ }
\newcommand{\cf}{cf.\ }
\newcommand{\resp}{respectively\ }
\newcommand{\Subsection}{\subsection}
\providecommand{\nobreakdash}{\nobreak\hbox{-}}
\setlength{\emergencystretch}{2em}
\multlinegap0pt
\newtheorem{lemma}{Lemma}
\title{Notes on Invariant Measures for Loop Groups}
\author{Doug Pickrell}
\date{Source: arXiv:2207.09913v6 (CC BY 4.0)}
\begin{document}
\maketitle

\noindent\emph{Source and licence.} Notes on Invariant Measures for Loop Groups. Version arXiv:2207.09913v6 (CC BY 4.0), distributed by arXiv under the Creative Commons Attribution 4.0 International licence, http://creativecommons.org/licenses/by/4.0/, as recorded in the arXiv licence metadata for this exact version and shown on the abstract page https://arxiv.org/abs/2207.09913v6. Full licence notice: \texttt{licenses/arxiv-2207.09913-CC-BY-4.0.txt}.\par\smallskip

\noindent\textit{Editorial scope.} Counted unit: the article's lemma lemma (source0.tex lines 2546--2550). The article's complete original proof of it is delivered with it (source0.tex lines 2553--2579, one \texttt{proof} environment of 27 lines). No mathematics is written, repaired or re-derived: the statement and the proof are the article's own lines, in the article's own notation, and no retained line is substituted or re-flowed.  No further source-local context is needed: every cross-reference of every retained line resolves inside the two delivered spans.  Nothing was omitted from any retained span, so the package carries no omission record.  No retained line contains a citation, so the wrapper carries no bibliography and no bibliography entry.  No retained line contains a figure environment, an image-inclusion command or drawing code, so no image is delivered and the package declares no asset dependency.  Editorial formatting only, in addition: an AMS \texttt{amsart} preamble that restates the article's own declarations and macros used by the retained lines, namely the environments \texttt{lemma} on separate counters, together with the standard packages those lines need.  The retained environments print in this document as Lemma 1 in order of appearance, whereas the article prints the same items with the numbering of its own full document.  The complete native source of the article as distributed in the arXiv e-print for this exact version is bundled unchanged as \texttt{sources/128476/source0.tex}.
\begin{lemma} Suppose $q<1$. Then the product
$$=\prod_{k=1}^{\infty}\left((\sum_{m,n=0}\mathbf m(m,n)
q^{m^2+m+m(n+\sqrt{n^2+16k^2})+(n+1)\frac{n+\sqrt{n^2+16k^2}}{2}-2k}\right)$$
converges and is nonzero.
\end{lemma}

\begin{proof} We must show that
$$\sum_{k=0}^{\infty}\left((\sum_{m>0,n=0}+2\sum_{m\ge 0,n\ge 1})
q^{m^2+m+m(n+\sqrt{n^2+16k^2})+(n+1)\frac{n+\sqrt{n^2+16k^2}}{2}-2k}\right)$$
converges. This can be split into two sums. The first sum is
$$\sum_{k=0}^{\infty}\sum_{m=0}^{\infty}
q^{m^2+m+4mk}$$
$$=\le \sum_{m=1}^{\infty}q^{m^2+m}(1-q^{4m})^{-1}$$
Since $(1-q^{4m})$ converges to $1$ as $m\uparrow \infty$, this sum is convergent.
The second sum is $2$ times
$$\sum_{k=0}^{\infty}\sum_{m\ge 0,n\ge 1}
q^{m^2+m+m(n+\sqrt{n^2+16k^2})+(n+1)\frac{n+\sqrt{n^2+16k^2}}{2}-2k}$$
$$\le \sum_{k=0}^{\infty}\sum_{m\ge 0,n\ge 1}
q^{m^2+m+m(n+4k)+(n+1)(n+2k)-2k}$$

$$\le \sum_{k=0}^{\infty}\sum_{m\ge 0,n\ge 1}
q^{m^2+m+m(n+4k)+(n+1)n+2nk}$$
$$=\sum_{m\ge 0,n\ge 1} \sum_{k=0}^{\infty}
q^{m^2+m+m(n+4k)+(n+1)n+2nk}$$
$$=\sum_{m\ge 0,n\ge 1}
q^{m^2+m+n^2+n+mn}(1-q^{4m+2k})^{-1}$$
Since $(1-q^{4m+2k})$ tends to $1$ as $m,n\uparrow \infty$, this will converge
iff
$$=\sum_{m\ge 0,n\ge 1}
q^{m^2+m+n^2+n+mn}$$
converges. This does converge, and this implies the lemma.

\end{proof}

\end{document}
